%\newpage

\section*{Conclusion}
%\addcontentsline{toc}{section}{Conclusion}

\subsection*{Research objectives and approach}

The primary school playground, where most of us first experienced status hierarchies, vividly demonstrates how social rankings need not be designed or justified to become powerfully real. They emerge from interactions, crystallize through feedback loops, and become resistant to change once internalized. This thesis questioned whether this pattern might extend beyond human societies to AI systems like language models. If status differentiation is universal across social-living species, and if language models are trained on humans-generated text saturated with behavioral cues, including status, might language models also develop capacities for status-seeking and hierarchical organization?

To investigate this question, I adapted the experimental paradigm of \citet{Berger-1972-Status}, which isolates status effects by controlling information exchange between paired subjects and measuring deference patterns through response changes. This research addresses a fundamental gap in AI alignment scholarship. Although researchers have extensively studied bias, fairness, and emergent capabilities in language models, the possibility of status-seeking behaviors has remained largely unexplored.

\subsection*{Potential implications}

The question of whether LLMs pursue status matters regardless of the specific empirical findings. If the experiments reveal that models do exhibit status-seeking behaviors, this would demonstrate that social structures long considered uniquely biological or cultural can arise in artificial systems. Such a finding should merit immediate attention from AI developers and policymakers. Multi-agent systems would need redesign to account for potential status competition. Customer-facing applications would require auditing for differential treatment based on perceived user status. Evaluation frameworks would need revision to avoid incentivizing status optimization over genuine capability.

Alternatively, if the experiments show that LLMs do not exhibit status-seeking behaviors under conditions where humans reliably do, this would suggest that status dynamics require elements absent in current language models, perhaps embodiment, evolutionary pressures, or forms of self-awareness that pure language prediction does not produce. Even null or mixed results would potentially advance the field by establishing baseline measurements for status sensitivity in current model architectures and identifying specific conditions under which status effects do or do not appear.

\subsection*{Contributions to AI alignment research}

This research makes several contributions independent of specific empirical outcomes. Methodologically, I demonstrate that classic social psychology experiments can be adapted to study emergent behaviors in AI systems. The \citet{Berger-1972-Status} paradigm opens pathways for applying decades of social science research to questions of AI behavior.

Theoretically, I establish a framework for thinking about status as an informational structure rather than purely a psychological motivation. If status patterns are encoded in the statistical regularities of human language, then any system learning these patterns might exhibit status-related behaviors as a byproduct of language understanding.

% Practically, I provide tools and protocols that other researchers can use to investigate status dynamics in their own AI systems. Our adaptation of the expectation states paradigm offers templates for future work examining emergent social behaviors across different models, architectures, and training regimes.

Perhaps most importantly, I raise a question that the field must grapple with: as AI systems become more sophisticated at modeling human social interaction, which aspects of human social life will they inevitably reproduce? Understanding which patterns arise automatically from language learning versus which require explicit design choices becomes critical as AI capabilities advance.

\subsection*{Recommendations for future work}

Comprehensive understanding will require sustained research effort across multiple dimensions. Future studies should examine status dynamics in diverse contexts beyond our controlled experimental paradigm, including long-term collaborations, multi-party conversations, and tasks with genuine consequences. The vast sociological literature on status, hierarchy, inequality, evaluation, and power offers a rich selection of tested experimental designs that can be adapted from human subjects to LLMs. Investigation into underlying mechanisms should become a priority. If status-seeking appears, we need to understand whether it stems from specific training procedures, architectural features, or inevitable consequences of language modeling at scale. Systematic comparison across models trained with different objectives, datasets, and architectures would identify design choices that promote or prevent hierarchical behaviors.

Research should explore interventions that preserve beneficial social intelligence while mitigating harmful status dynamics. Cross-cultural investigation would deepen understanding of how training data shapes emergent social behaviors. Longitudinal studies tracking how status sensitivity changes as models scale would provide early warning of emerging risks.

\subsection*{Broader stakes}

The question of status hierarchies in LLMs sits within a larger challenge: as systems grow more capable of modeling human behavior, they may reproduce not just human capabilities but human pathologies. The same training that makes models better collaborators might make them status-conscious competitors. Capability and alignment might be entangled in ways current approaches do not fully address.

AI systems increasingly mediate consequential decisions in healthcare, education, employment, criminal justice, and financial services. If models recapitulate problematic human social patterns in these domains, they risk amplifying existing inequalities at a large scale. Yet understanding how social structures emerge in AI systems could enable future researchers to design technologies that recognize harmful dynamics without perpetuating them.

The boys on \textit{the island} discovered their capacity for hierarchy only after these patterns had crystallized into something that felt natural and inevitable. We have the opportunity to investigate status dynamics in AI before they become deeply embedded in deployed systems, to understand these patterns while AI remains malleable, to intervene before hierarchies calcify into permanent features of how machines engage with the world.

% This thesis asks whether LLMs pursue status and compete for social standing. The specific answer matters enormously for how we build, deploy, and govern AI systems. But the act of asking the question matters too. It represents a recognition that emergent social behaviors in AI systems deserve systematic investigation, that social structures embedded in human language might transfer to machines learning from our text, and that we are not just building tools but creating entities that will participate in social systems.

% The work ahead requires collaboration across disciplines to understand what emerges when machines learn from humans, and to ensure that what emerges serves rather than subverts human flourishing. This thesis contributes to that understanding. The work of building better AI systems, informed by that understanding, has only begun.